\section{System Model}
\label{sec:system-model}

\subsection{Enterprise Chain Architecture}

The Basis Network zkEVM L2 follows a hub-and-spoke topology (established
in RU-L7~\cite{rul7-hub-spoke}) where each enterprise operates an
independent chain. The settlement layer (Basis Network L1, an Avalanche
Subnet-EVM instance) serves as the hub, and each enterprise chain is a
spoke. Each spoke:

\begin{enumerate}
  \item Sequences transactions from its enterprise's applications
    (e.g., PLASMA industrial maintenance, Trace ERP).
  \item Maintains an independent Sparse Merkle Tree of enterprise state
    (depth 32, Poseidon hash, as established in RU-L1).
  \item Produces halo2-KZG validity proofs attesting to correct state
    transitions (proof system selected in RU-L9~\cite{rul9-plonk}).
  \item Submits proofs and state roots to the L1 for verification.
\end{enumerate}

\subsection{Verification Cost Model}

Let $G(s)$ denote the L1 verification gas cost for proof system $s$.
Based on published benchmarks and our prior measurements:

\begin{align}
  G(\text{halo2-KZG}) &= 420\text{K gas} \label{eq:halo2-gas} \\
  G(\text{Groth16}) &= 207\text{K} + 7.16\text{K} \cdot l \label{eq:groth16-gas} \\
  G(\text{PLONK-KZG}) &= 290\text{K gas} \label{eq:plonk-gas}
\end{align}

\noindent where $l$ is the number of public inputs. For $N$ enterprises
submitting individually, total verification gas is:

\begin{equation}
  G_{\text{total}} = N \cdot G(s) \label{eq:total-gas}
\end{equation}

\noindent At $N{=}8$ with halo2-KZG: $G_{\text{total}} = 8 \times 420\text{K} = 3.36\text{M gas}$.

\subsection{Aggregation Model}

An aggregation strategy $\mathcal{A}$ takes $N$ proofs
$\{\pi_1, \ldots, \pi_N\}$ and produces a single aggregated proof
$\pi_{\text{agg}}$ verifiable on L1. We define:

\begin{itemize}
  \item $G_{\mathcal{A}}$: L1 gas cost to verify $\pi_{\text{agg}}$
    (ideally independent of $N$).
  \item $T_{\mathcal{A}}(N)$: Wall-clock time to produce
    $\pi_{\text{agg}}$ from $N$ inner proofs.
  \item $|\pi_{\mathcal{A}}|$: Size of $\pi_{\text{agg}}$ in bytes.
  \item $M_{\mathcal{A}}(N)$: Peak memory during aggregation.
\end{itemize}

The gas savings factor is:

\begin{equation}
  S(N) = \frac{N \cdot G(s)}{G_{\mathcal{A}}} \label{eq:savings}
\end{equation}

\noindent and the amortized per-enterprise gas cost is:

\begin{equation}
  G_{\text{amort}}(N) = \frac{G_{\mathcal{A}}}{N} \label{eq:amort}
\end{equation}

\subsection{Threat Model}

The aggregation layer introduces a new trust assumption: the aggregator
must faithfully include all valid enterprise proofs. We assume:

\begin{itemize}
  \item \textbf{Honest aggregator}: The aggregator does not forge proofs
    or selectively exclude valid proofs. Censorship resistance is
    addressed by forced inclusion mechanisms (out of scope for this
    paper).
  \item \textbf{Adversarial enterprises}: Individual enterprises may
    submit invalid proofs. The aggregation must reject these without
    invalidating other enterprises' proofs.
  \item \textbf{Trusted L1 verifier}: The L1 smart contract correctly
    implements the verification algorithm. The Groth16 verifier is
    audited; the halo2-KZG verifier follows production patterns from
    Axiom and Scroll.
\end{itemize}

\subsection{Success Criteria}

Based on the system model and operational requirements, we define the
following success criteria:

\begin{enumerate}
  \item Gas reduction $\geq 4\times$ for $N{=}8$ enterprises compared
    to individual verification.
  \item Aggregation time $< 30$\,s for $N{=}8$.
  \item Final proof size $< 2$\,KB regardless of $N$.
  \item Soundness: the aggregated proof is valid if and only if all
    component proofs are valid.
\end{enumerate}
